% [4 pts] ~2.6 páginas. Identificación y CUANTIFICACIÓN de riesgos con cita a Nota/página.
% Estructura por riesgo: fuente (qué partida) → dirección (qué movimiento daña) → magnitud → coberturas naturales.
\section{Identificación de riesgos financieros}

La Nota 3.1 clasifica la exposición de PETROPERÚ en riesgos de mercado (tipo de cambio, tasa de interés y precio
del crudo), de crédito y de liquidez. El Cuadro~\ref{tab:matriz} resume su magnitud y si pueden gestionarse
eficientemente con derivados.

\begin{table}[htbp]\centering\small
\caption{Matriz de riesgos financieros de PETROPERÚ al cierre de 2025}\label{tab:matriz}
\newcommand{\muestra}[1]{\fcolorbox{gray!60}{#1}{\rule{0pt}{0.55em}\rule{0.9em}{0pt}}}
\renewcommand{\arraystretch}{1.25}
\setlength{\fboxsep}{0pt}
\newcommand{\sep}{\arrayrulecolor{white}\specialrule{1.2pt}{0pt}{0pt}\arrayrulecolor{black}}
\begin{tabularx}{\linewidth}{>{\raggedright\arraybackslash\bfseries}p{1.85cm}
    >{\raggedright\arraybackslash}X
    >{\raggedright\arraybackslash}p{3.0cm}
    >{\raggedright\arraybackslash}p{3.45cm}
    >{\raggedright\arraybackslash}p{3.0cm}}
\toprule
\normalfont\textbf{Riesgo} & \textbf{Origen} & \textbf{Exposición} & \textbf{Impacto (US\$ MM)} & \textbf{Respuesta} \\
\midrule
\rowcolor{prioAlta}
Tipo de cambio & Pasivo neto en S/: préstamo del Banco de la Nación (BN) y cuentas por pagar
  & S/~\PenNetoAbs\ MM \newline (US\$~\PenNetoUSD\ MM)
  & TC $-10$\,\%: (\SensApreciacion) \newline TC $+10$\,\%: \SensDepreciacion
  & CCS + NDF \newline (Estrategia 1) \\ \sep
\rowcolor{prioAlta}
Precio del crudo & Inventario de crudo; compras de \ComprasMBDC\ MBDC
  & \InvCrudoMbl\ MBL \newline (US\$~\InvCrudoUSD\ MM)
  & WTI $-20$\,\%: (\PerdidaInvVeinte) \newline = \PerdidaInvVeinteXUB\ veces la utilidad bruta
  & Swap o collar OTC \newline (Estrategia 2) \\ \sep
\rowcolor{prioMedia}
Tasa de interés & Bonos a tasa fija; refinanciación del CESCE y de líneas de corto plazo
  & Bonos: US\$~\BonosVR\ MM a mercado \newline (US\$~\BonosLibros\ MM en libros)
  & $\pm100$~pb: $\mp$\BonosDVCien\ en el valor razonable
  & Solo la tasa base, no el spread \\ \sep
\rowcolor{prioNo}
Liquidez & Capital de trabajo negativo; líneas usadas al \LineasPct
  & US\$~\VencMenosUnAnio\ MM de pagos a menos de un año
  & Restricción, no riesgo de precio
  & Condiciona el diseño: sin márgenes ni primas altas \\ \sep
\rowcolor{prioNo}
Crédito & Mayoristas y entidades del Estado
  & Baja concentración (Nota 3.1(b))
  & No significativo
  & No es necesario \\
\bottomrule
\end{tabularx}
\par\vspace{3pt}{\footnotesize\raggedright\textbf{Prioridad:}\quad
  \muestra{prioAlta}~Alta: se cubre con derivados\qquad
  \muestra{prioMedia}~Media: cobertura parcial\qquad
  \muestra{prioNo}~No se cubre con derivados\par}
\fuente{\textcite{petroperu2026eeff}, Notas 3, 10 y 14 (pp.~42--46, 60--61 y 74--78). Elaboración propia.}
\end{table}

\subsection{Riesgo de tipo de cambio}
Con contabilidad en dólares, el riesgo surge de las partidas en soles, euros y yenes. La exposición relevante es en
soles: una posición pasiva neta de S/~\PenNetoAbs\ millones (US\$~\PenNetoUSD\ millones al TC SBS de \TCcierre\ que usa el EEFF),
\PenNetoCrec\ mayor que la de 2024 (S/~\PenNetoVeinticuatro\ millones). La concentra el rubro \emph{otros pasivos
financieros} (S/~\PenOtrosPasFin\ millones), que incluye el préstamo del BN con garantía soberana por
S/~\BNSaldoPEN\ millones al \BNTasa, pagadero en \BNCuotasTot\ cuotas mensuales de enero de 2025 a diciembre de 2028
(Nota~14(ii), p.~74); como en 2025 solo se pagaron intereses, la sección~3 supone \BNCuotasSup\ cuotas desde 2026. Las posiciones en euros (activa por EUR~\EurNeto\ millones) y yenes (pasiva por
JPY~\JpyNeto\ millones) no son significativas según la Gerencia (Figura~\ref{fig:posicion}).

\begin{figure}[htbp]\centering
  \includegraphics[width=0.86\linewidth]{f_posicion_pen}
  \caption{Composición de la posición en soles al cierre de 2025 (S/ millones)}\label{fig:posicion}
  \fuente{\textcite{petroperu2026eeff}, Nota 3.1(a)(i), p.~42. Elaboración propia.}
\end{figure}

El riesgo ya se materializó: en 2025 el tipo de cambio interbancario medio del BCRP pasó de \TCcierreBCRPVeinticuatro\ a
\TCcierreBCRPVeinticinco\ (\TCvarVeinticinco; el sol se apreció, Figura~\ref{fig:historia}) y la diferencia de
cambio aumentó en US\$~\DifCambioBN\ millones el saldo en dólares del préstamo del BN (Nota~14(c), p.~77), aunque en
el agregado la Compañía reportó una ganancia cambiaria neta de US\$~\GananciaDifCambio\ millones. La Compañía estima
que un movimiento de $\pm$10\,\% del dólar frente al sol afecta el resultado antes de impuestos en
$\pm$US\$~\SensibilidadFX\ millones, una aproximación lineal. El cálculo exacto es asimétrico: $-$US\$~\SensApreciacion\
millones si el TC cae 10\,\% y +US\$~\SensDepreciacion\ millones si sube 10\,\%. El descalce no es solo contable: las ventas
locales se facturan en soles, pero sus precios siguen la paridad de importación en dólares (Nota~1-c), por lo que
económicamente son flujos en dólares que no cubren de forma natural una deuda en soles a tasa fija.

\begin{figure}[htbp]\centering
  \includegraphics[width=0.9\linewidth]{f_historia}
  \caption{Tipo de cambio y precio del WTI, promedios mensuales (\HistDesde\ a \HistHasta)}\label{fig:historia}
  \fuente{\textcite{bcrpseries2026}, series PN01210PM y PN01660XM. Volatilidad anualizada de los retornos
  logarítmicos mensuales: TC \VolHistTC; WTI \VolHistWTI. Elaboración propia.}
\end{figure}

\subsection{Riesgo de precio del crudo}
Es el riesgo operativo-financiero más directo. En 2025 PETROPERÚ compró crudo y productos por
US\$~\ComprasUSD\ millones, a un precio promedio de US\$~\PrecioCompras/bl y con un volumen de \ComprasMBDC\ MBDC
(Nota~23, p.~93), pero el crudo cerró el año en US\$~\WTIcierreVeinticinco/bl, \WTIvarVeinticinco\ frente a los
US\$~\WTIcierreVeinticuatro/bl de 2024 (Nota~10, p.~61). Entre la compra y la venta del producto refinado pasan
varias semanas, y en ese lapso el inventario pierde valor si el precio cae; la Nota~1-d señala esta caída como una
causa de la pérdida de 2025. Al cierre, el inventario de crudo era de \InvCrudoMbl\ MBL (US\$~\InvCrudoUSD\ millones)
y el de hidrocarburos sumaba US\$~\InvHidroUSD\ millones. Con una utilidad bruta de apenas
US\$~\UtilidadBruta\ millones, una caída de 20\,\% del WTI sobre el inventario de crudo (US\$~\PerdidaInvVeinte\
millones) equivale a \PerdidaInvVeinteXUB\ veces la utilidad bruta del año. El riesgo es simétrico y de gran escala:
en 2026 el WTI subió hasta US\$~\WTISpot/bl al \FechaValorizacion, en plena crisis de oferta en Medio Oriente
\parencite{reuters2026oil}. El FEPC solo alcanza al Diésel BX y al Petróleo Industrial~6 y en 2025 fue un aporte
neto (\FEPCpct\ de los ingresos), no una compensación (Nota~1-c, p.~19).

\subsection{Riesgos de tasa de interés, liquidez y crédito}
Toda la deuda está pactada a tasa fija (Figura~\ref{fig:deuda}): bonos 2032 (\CuponTreintaDos) y 2047 (\CuponCuarentaSiete) por
US\$~\BonosNominal\ millones nominales, el préstamo CESCE (\CESCETasa, 2030), el préstamo del BN y líneas de corto
plazo. No hay riesgo de flujo de caja por tasas variables; el riesgo está en el valor razonable y en la
refinanciación. Los bonos valen US\$~\BonosVR\ millones a mercado frente a US\$~\BonosLibros\ millones en libros
(Nota~14(d), p.~78), lo que implica un rendimiento de \YTMBonos\ (cálculo propio): un spread de crédito de unos
\SpreadBonosPb~pb sobre el Tesoro de EE.\,UU. de igual duración (\USTDurDic). Con una duración modificada de \BonosDurMod\ años, cada 100~pb modifican su valor en unos
US\$~\BonosDVCien\ millones. Además, por la opinión con salvedades se incumplió un covenant no financiero del CESCE y su
saldo (US\$~\CESCESaldo\ millones) se reclasificó al corto plazo. Un swap puede fijar la tasa base, pero no el spread,
que es el componente dominante.

\begin{figure}[htbp]\centering
  \includegraphics[width=0.86\linewidth]{f_deuda}
  \caption{Deuda financiera y con el accionista al cierre de 2025 (US\$~\DeudaTotal\ millones, sin intereses devengados)}\label{fig:deuda}
  \fuente{\textcite{petroperu2026eeff}, Nota 14(a), p.~76, y Nota 3.2, p.~46. Elaboración propia.}
\end{figure}

El riesgo de liquidez es crítico: de US\$~\LineasRev\ millones en líneas revolventes se usaban
US\$~\LineasUsadas\ millones (\LineasPct), y los pagos no descontados a menos de un año suman
US\$~\VencMenosUnAnio\ millones (Nota 3.1(c), pp.~44--45). No se cubre con derivados, pero condiciona su diseño:
deben evitarse instrumentos con llamadas de margen diarias, como los futuros en bolsa, y primas iniciales altas.
El riesgo de crédito es bajo, porque la cartera se concentra en mayoristas de primer nivel y entidades del Estado.
